%% GENERATED FILE -- do not hand-edit.
%% SOURCE: experiments/019-perturb-seq-costing/scripts/render_tex_tables.py
\begin{table}[htbp]\centering
\footnotesize
\caption[]{Measured response of the transcriptome to a gene deletion. A single deletion moves a few hundred genes by 1.25$\times$ but only $\sim$10--15 by 2$\times$, so designing at two-fold targets a few percent of the response. The median responding gene moves $\sim$1.34$\times$, and because the power coefficient $A$ scales as $|\Delta|^{-2}$ this multiplies every cell requirement by $\sim$5.6 relative to the two-fold assumption. Doubling the perturbation count roughly doubles the responder count at 1.25$\times$ and trebles it at 2$\times$, without shifting the median response --- more perturbations make a cell noisier, not its individual effects larger. Generated from \file{experiments/019-perturb-seq-costing/scripts/effect_size_analysis.py}.}\label{tab:effect-size}
\begin{tabular}{lrrrrrr}
\toprule
Compendium & Strains & \multicolumn{3}{c}{Median responding genes at} & \multicolumn{2}{c}{Median $|\Delta|$, responders} \\
 & & 1.25$\times$ & 1.5$\times$ & 2$\times$ & log$_2$ & fold \\
\midrule
Kemmeren 2014, single deletion & 1{,}484 & 269 & 55 & 15 & \textbf{0.42} & 1.34$\times$ \\
Sameith 2015, single deletion & 82 & 220 & 42 & 10 & \textbf{0.42} & 1.34$\times$ \\
Sameith 2015, \textbf{double} deletion & 72 & 460 & 112 & 30 & \textbf{0.44} & 1.36$\times$ \\
\bottomrule
\end{tabular}
\end{table}
